\documentclass{article}
\usepackage[T1]{fontenc}
\usepackage{lmodern}
\usepackage{graphicx}
\usepackage{amsmath,amssymb}
\usepackage[a4paper,margin=2cm]{geometry}
\usepackage[absolute,overlay]{textpos}
\setlength{\TPHorizModule}{1mm}
\setlength{\TPVertModule}{1mm}
\usepackage[indonesian]{babel}
\usepackage{hyperref}
\setlength{\emergencystretch}{2em}
% unit: Halaman 27

\begin{document}

% === Halaman 27 (python/native) ===

5. Replay Attack

• Tujuan: menggunakan kembali pesan sah yang pernah dikirim sebelumnya untuk 

menyebabkan suatu tindakan terjadi lagi.

• Yang menarik, pada replay attack penyerang tidak perlu memecahkan enkripsi.

• Untuk mencegah replay attack, biasanya digunakan mekanisme yang membuat 

setiap pesan mempunyai keunikan atau urutan, misalnya: Nonce, Timestamp, 

atau Sequence number 

• Misalnya pesan diberi sequence number: 

      Message 101

      Message 102

      Message 103

Jika penyerang mengirim kembali Message 101 setelah sistem sudah menerima 

Message 103, sistem dapat menolaknya karena nomor tersebut sudah 

kedaluwarsa.

27

\end{document}
